% Section 3.1, from lectures/L03/appendix/a_filters.md.

\section{Passive filters, cascading, and Q}
\label{c3:app:a}

The first useful circuits in the book, and the third appearance of the one piece of arithmetic that decides everything.

\subsection{Two components, four filters}
\label{c3:sec:a1}
A resistor and a capacitor in series make a divider whose ratio depends on frequency. Which component you take the output across decides the kind of filter.

\[
  H_{LP} = \frac{1}{1 + j f/f_c}, \qquad H_{HP} = \frac{j f/f_c}{1 + j f/f_c}
\]

Both have the same corner, $f_c = 1/(2\pi RC)$, and both are 3 dB down there. At that corner the low-pass lags 45 degrees and the high-pass leads 45, so \textbf{their outputs add to give the input back exactly, at every frequency.} That is the cleanest statement of what complementary means.

\bookfigure{lectures/L03/appendix/images/filter_family.png}{Low-pass, high-pass and band-pass magnitude responses against frequency normalised to a common corner, on logarithmic axes. The low-pass is flat then falls, the high-pass rises then is flat, and the band-pass rises and falls, reaching nearly zero decibels between two corners two decades apart.}{c3:fig:filterfamily}

A \textbf{band-pass} is a high-pass and a low-pass in series. It reaches 0 dB in the middle only if the corners are well apart: a decade apart falls 0.83 dB short, equal corners 6 dB.

\subsection{A filter's own impedances}
\label{c3:sec:a2}
A filter is a two-port and both of its ports have a frequency-dependent impedance. Both matter for Chapter~\ref{c1:ch:circuits}'s reason: the source upstream and the load downstream form dividers with them.

\[
  Z_{in} = R + \frac{1}{j\omega C}, \qquad Z_{out} = R \parallel \frac{1}{j\omega C}
\]

$Z_{in}$ is capacitor-dominated and large at low frequency, approaching $R$ at high; $Z_{out}$ is $R$ at low frequency and approaches zero at high. \textbf{So a low-pass RC presents its own series resistance as an output impedance at low frequency, and the next stage sees it.}

\subsection{Cascading, and the corner you did not design}
\label{c3:sec:a3}
Two identical low-pass sections in a row look like they should give the square of one section's response: 6 dB down at the corner, falling at 40 dB per decade. \textbf{That is the answer for two sections with a buffer between them, and it is wrong for two connected directly.}

\bookfigure{lectures/L03/appendix/images/cascade_loading.png}{Magnitude against frequency for one RC section, for two sections with a buffer between them, and for two sections cascaded directly. The directly cascaded pair falls away earliest, and a marker shows it is three decibels down at 0.37 of one section's corner.}{c3:fig:cascadeloading}

The second section loads the first. Its input impedance parallels the first section's capacitor, so the two poles split apart, their product staying at $f_c^2$:

\[
  f_{low} = f_c \frac{3 - \sqrt{5}}{2} = 0.382 f_c, \qquad f_{high} = f_c \frac{3 + \sqrt{5}}{2} = 2.618 f_c
\]

\begin{booktable}{@{}Ll@{}}
  \thead{Arrangement} & \thead{3 dB point} \\
  \midrule
  One section & $1.000 f_c$ \\
  Two sections, buffered & $0.644 f_c$ \\
  Two sections, cascaded directly & $0.374 f_c$ \\
\end{booktable}

The direct cascade is a factor of \textbf{1.72} lower than the buffered pair. \textbf{A filter designed by multiplying two responses together and built by soldering two sections together is not the filter that was designed.}

\textbf{The fix is a buffer:} high input impedance, low output impedance, between the sections so the second cannot load the first. An op-amp follower is exactly that, which is why this chapter covers filters and op-amps together.

\textbf{This is Chapter~\ref{c1:ch:circuits}'s arithmetic for the third time.} A divider loaded by 10 kilohm loses a third of its output; here two filter sections move each other's poles. In Chapter~\ref{c10:ch:amplifier} an op-amp loses 11 dB the same way.

\subsection{LC, resonance, and Q}
\label{c3:sec:a4}
An inductor and a capacitor have reactances of opposite sign, so at one frequency they cancel:

\[
  f_0 = \frac{1}{2\pi\sqrt{LC}}
\]

In series they cancel to zero impedance, in parallel to infinite. \textbf{Nothing about that frequency depends on any resistance,} which is what makes a resonance different in kind from an RC corner. The resistance decides only the sharpness, for a series RLC taken across the resistor:

\[
  Q = \frac{1}{R}\sqrt{\frac{L}{C}}, \qquad \text{bandwidth} = \frac{f_0}{Q}
\]

\bookfigure{lectures/L03/appendix/images/resonance_q.png}{Band-pass magnitude for three values of Q on the same resonance, showing a broad peak at Q of one, a narrower one at Q of 3.3 and a sharp one at Q of 10, all reaching zero decibels at resonance.}{c3:fig:resonanceq}

\textbf{Q is two things at once, and the second is the one that gets people.} It is the sharpness of the peak, and it is the factor by which the inductor's and the capacitor's own voltages exceed the input: at resonance the LC is a short, so the current $V/R$ flows through a reactance $Q$ times $R$.

\[
  V_L = Q \times V_{in}
\]

\textbf{A Q of 10 with 10 V applied puts 100 V across a capacitor.} That is how a filter correct on paper destroys a component on a bench, and it is nowhere in the transfer function, which describes only the output.

\subsection{Higher orders, and what this book does not do}
\label{c3:sec:a5}
Two poles is as far as this book goes, and it gets there by cascading rather than by designing.

Real design starts from a specification, picks the polynomial that meets it with the fewest poles, and realises it as a circuit: Butterworth for a flat passband, Chebyshev for a steeper transition at the cost of ripple, Bessel for flat group delay. Chapter~\ref{c4:ch:feedback} builds one Sallen-Key, the cheapest complex pole pair from one amplifier. Left out: pole placement, the approximation problem, tolerance sensitivity, switched-capacitor realisations.

\subsection{What this section is blind to}
\label{c3:sec:a6}
\begin{itemize}
  \item \textbf{Component tolerance.} Every corner here is nominal. A filter of 5 per cent parts has a corner good to about 7 per cent, and a high-Q filter is far worse, because Q depends on a ratio of two square roots.
  \item \textbf{Real inductors.} The Q values assume the resistance is the one you put there. A real inductor's winding resistance is often dominant, and it caps Q where no external resistor can lift it.
  \item \textbf{Transients.} A high-Q filter rings, for roughly $Q$ cycles. This book never computes it.
\end{itemize}
